\section{The canvas}
\label{sec:canvas}

The canvas fills the centre of the main window (callout~3 in Figure~\ref{fig:main-window}). It
shows the current frame: before a run, the image with its mask as a blue tint; while you edit a
mask, the same with the editing banner (Section~\ref{sec:roi-editing}); after a run, the
displacement field over the image, with a colorbar (Section~\ref{sec:results}). A toolbar runs along its
top edge and the frame navigator along its bottom edge.

\subsection{Zoom and pan}

\begin{description}[leftmargin=1.2cm,style=nextline]
  \item[Fit]
    Fits the whole image into the canvas. RAFTcorr also fits the image when you open a folder.
  \item[100\%]
    Shows one image pixel per screen pixel.
  \item[Zoom in, Zoom out]
    The two magnifier buttons next to \ui{100\%} zoom by one step, a factor of 1.15.
  \item[Mouse wheel]
    Zooms by the same step around the point under the cursor, which stays in place.
  \item[Middle mouse button]
    Drag with the middle button (the pressed wheel) to pan. The canvas has no scroll bars.
\end{description}
The wheel and the zoom buttons stop at 10\,\% and 2000\,\%. Without a middle button, as on a touchpad, move
around by zooming out (or \ui{Fit}) and zooming in again over the place you want to see.

At the right end of the toolbar, a read-out follows the cursor: its position as
\file{x 512.0  y 384.0 px}, and, when a field is shown, the value of the field under the cursor
with its unit. The position is in image coordinates, $x$ to the right and $y$ downward in pixels,
with the centre of the top-left pixel at $(0, 0)$, the same convention as the data files.

\subsection{The configuration card}

Once at least two images are loaded, a card in the top-left corner of the canvas summarises what
a run would do:
\begin{description}[leftmargin=1.2cm,style=nextline]
  \item[MODE] \ui{Accumulative}, \ui{Incremental}, \ui{Incremental, every N frames} (with the
    interval) or \ui{Custom reference frames} (Section~\ref{sec:workflow}).
  \item[MODEL] The selected model, for example \ui{RAFTcorr-s4}.
  \item[ITERATIONS] The number of iterations (Section~\ref{sec:parameters}).
\end{description}
The card follows the sidebar as you change the settings. Mouse clicks pass through it to the
canvas.

\subsection{The frame navigator}

The bar below the canvas holds, from left to right:
\begin{description}[leftmargin=1.2cm,style=nextline]
  \item[Previous frame, Play animation, Next frame]
    The arrows step one frame back or forward and stop at the first and the last frame. Play
    advances one frame at a time and starts again at frame~1 after the last frame; the same
    button pauses. Opening other images stops the playback.
  \item[Playback speed]
    \file{1 fps}, \file{2 fps} (default), \file{5 fps}, \file{10 fps} or \file{30 fps}. This is
    the speed of the display, not the frame rate of the experiment; that one belongs to the
    physical units (Section~\ref{sec:physical-units}).
  \item[FRAME n/N]
    The current frame and the number of frames, counted from 1 as in the image list.
  \item[Slider]
    Drag it to move through the sequence.
\end{description}
The slider, the image list and the canvas always show the same frame: moving one moves the
others.

RAFTcorr has no keyboard shortcut for frames that works everywhere in the window. \kbd{Up},
\kbd{Down}, \kbd{Page Up}, \kbd{Page Down}, \kbd{Home} and \kbd{End} move through the frames when
the image list has the keyboard focus (click a row first), and the same keys move the slider after
you click it. Table~\ref{tab:shortcuts} lists every key the application answers to.
% CHECK: slider keys are Qt's standard QSlider behaviour (singleStep 1, pageStep 10 by default), not RAFTcorr code; confirm on Windows and macOS
